\providecommand{\Needspace}[1]{}
\Needspace{8\baselineskip}
\item \textbf{Ordenamiento y busqueda de periodos orbitales}.
\nopagebreak[4]

\begin{tcolorbox}[colback=blue!2!white,colframe=blue!55!black,title={Enunciado},
                  before skip=3pt,after skip=5pt,boxsep=3pt,top=3pt,bottom=3pt,
                  breakable]
\small
\raggedright
Un equipo observa exoplanetas y registra sus periodos orbitales en dias:
\[
\texttt{periodos = [12.4,\ 3.1,\ 8.7,\ 25.0,\ 1.9,\ 14.2,\ 6.5,\ 18.3,\ 30.5,\ 4.8]}.
\]

Utilice los algoritmos entregados a continuacion. No debe escribir otro
algoritmo de ordenamiento ni otro algoritmo de busqueda:

\begin{pythonjupyter}
def bubble_sort(lista):
    n = len(lista)
    for i in range(n - 1):
        for j in range(n - 1 - i):
            if lista[j] > lista[j + 1]:
                lista[j], lista[j + 1] = lista[j + 1], lista[j]
    return lista

def binary_search(lista_ordenada, objetivo):
    low, high = 0, len(lista_ordenada) - 1

    while low <= high:
        mid = (low + high) // 2

        if lista_ordenada[mid] == objetivo:
            return mid
        elif objetivo < lista_ordenada[mid]:
            high = mid - 1
        else:
            low = mid + 1

    return -1
\end{pythonjupyter}

Realice las siguientes operaciones:

\begin{itemize}
\setlength{\topsep}{4pt}
\setlength{\itemsep}{3pt}
\setlength{\parsep}{0pt}
\item[(a)] Defina la lista \texttt{periodos}.
\item[(b)] Ordene los periodos de menor a mayor usando \texttt{bubble\_sort}.
Debe ordenar una copia de la lista original, no la lista original directamente.
\item[(c)] Imprima la lista original y la lista ordenada. Verifique con un
\texttt{print} que la lista original no cambio.
\item[(d)] Usando solamente la lista ordenada, determine el periodo minimo, el
periodo maximo y la mediana. Como la lista tiene una cantidad par de datos, la
mediana corresponde al promedio de los dos valores centrales.
\item[(e)] Use \texttt{binary\_search} para buscar el periodo \texttt{14.2} en
la lista ordenada. Imprima su posicion si fue encontrado.
\item[(f)] Use \texttt{binary\_search} para buscar el periodo \texttt{10.0}.
Si no aparece, imprima un mensaje claro indicando que no fue encontrado.
\item[(g)] Recorra la lista ordenada y construya una nueva lista llamada
\texttt{ventana\_habitable} con los periodos entre \texttt{5.0} y \texttt{20.0}
dias, incluyendo los extremos si aparecen.
\item[(h)] Imprima la lista \texttt{ventana\_habitable}, su cantidad de
elementos y una frase que explique por que \texttt{binary\_search} solo tiene
sentido despues de ordenar.
\end{itemize}
\end{tcolorbox}

\begin{solution}
\begin{pythonjupyter}
periodos = [12.4, 3.1, 8.7, 25.0, 1.9,
            14.2, 6.5, 18.3, 30.5, 4.8]

def bubble_sort(lista):
    n = len(lista)
    for i in range(n - 1):
        for j in range(n - 1 - i):
            if lista[j] > lista[j + 1]:
                lista[j], lista[j + 1] = lista[j + 1], lista[j]
    return lista

def binary_search(lista_ordenada, objetivo):
    low, high = 0, len(lista_ordenada) - 1

    while low <= high:
        mid = (low + high) // 2

        if lista_ordenada[mid] == objetivo:
            return mid
        elif objetivo < lista_ordenada[mid]:
            high = mid - 1
        else:
            low = mid + 1

    return -1

periodos_originales = periodos.copy()
periodos_ordenados = bubble_sort(periodos.copy())

print("Lista original:", periodos)
print("Lista ordenada:", periodos_ordenados)
print("La original no cambio:", periodos == periodos_originales)

periodo_minimo = periodos_ordenados[0]
periodo_maximo = periodos_ordenados[-1]

n = len(periodos_ordenados)
mediana = (
    periodos_ordenados[n//2 - 1] +
    periodos_ordenados[n//2]
) / 2

print("Periodo minimo:", periodo_minimo)
print("Periodo maximo:", periodo_maximo)
print("Mediana:", mediana)

posicion_142 = binary_search(periodos_ordenados, 14.2)
if posicion_142 == -1:
    print("El periodo 14.2 no fue encontrado")
else:
    print("Posicion de 14.2:", posicion_142)

posicion_100 = binary_search(periodos_ordenados, 10.0)
if posicion_100 == -1:
    print("El periodo 10.0 no fue encontrado")
else:
    print("Posicion de 10.0:", posicion_100)

ventana_habitable = []
for p in periodos_ordenados:
    if 5.0 <= p <= 20.0:
        ventana_habitable.append(p)

print("Periodos en ventana:", ventana_habitable)
print("Cantidad:", len(ventana_habitable))
print("La busqueda binaria funciona sobre una lista ordenada porque")
print("cada comparacion con el valor central permite descartar una mitad.")
print("Sin orden, descartar una mitad no estaria justificado.")
\end{pythonjupyter}
\end{solution}
